\section{Introduction}
\label{sec:intro}

Volumetric medical image segmentation relies on decoders to fuse multi-scale
features and restore spatial detail.  This restoration is expensive: every
high-resolution 3D convolution operates over millions of voxels, often with
wide feature maps.  Modern architectures therefore face a tension between
anatomical accuracy and the memory, latency, and energy required for deployment.
This tension has motivated efficient architectures and automated configuration
pipelines such as nnU-Net~\cite{isensee2021nnunet}, as well as transformer and
large-kernel models including UNETR~\cite{hatamizadeh2022unetr}, Swin
UNETR~\cite{tang2022swinunetr}, and 3D UX-Net~\cite{lee20233duxnet}.

EffiDec3D~\cite{rahman2025effidec3d} recently showed that decoder channels and
high-resolution stages can be aggressively reduced while retaining competitive
aggregate Dice.  Its 3D UX-Net variant reports a 96.4\% reduction in decoder
parameters and a 93.0\% reduction in decoder FLOPs.  This is strong evidence
that conventional decoders are over-provisioned in aggregate.  Aggregate
metrics, however, cannot reveal whether the remaining utility is uniformly
distributed.  A static decoder assigns the same architecture to every subject,
organ, and location, although the need for detail restoration may vary sharply.

This distinction matters.  A difficult or uncertain voxel is not necessarily a
voxel that benefits from a stronger decoder: both models may fail there.
Conversely, model disagreement is not necessarily improvement, because the full
decoder can destroy a prediction that the efficient decoder got right.  We
therefore ask:
\begin{quote}
\emph{Where is the marginal utility of additional decoder capacity, how
concentrated is it, and can it be identified without access to ground truth?}
\end{quote}

We answer this question with a controlled marginal-utility study.  We compare a
full 3D UX-Net decoder (\fullmodel) with an EffiDec3D-compressed decoder
(\effimodel) under a common dataset, split, preprocessing pipeline, training
schedule, and evaluation protocol.  At each voxel we distinguish positive,
negative, and net transitions between their predictions.  We then analyze this
utility by anatomy, boundary proximity, uncertainty, training stage, and
selection budget.  Statistical comparisons are performed at the subject level
to avoid treating millions of correlated voxels as independent samples.

Our current contributions are:
\begin{itemize}
    \item We formulate \emph{decoder marginal utility}, separating additional
    decoder capacity from generic segmentation error and explicitly accounting
    for both corrected and damaged predictions.
    \item We provide a controlled audit of full and efficient 3D UX-Net decoders
    across two efficient-decoder seeds, reporting accuracy, HD95, parameters,
    MACs, latency, and memory together with the observed reproduction gap.
    \item We characterize the anatomical and temporal distribution of decoder
    opportunity, including boundary/interior error, organ size, uncertainty,
    and training-checkpoint evolution.
    \item We define a selective-allocation protocol with predeclared budgets,
    paired subject bootstrap, net utility, and contiguous 3D regions.  This
    protocol distinguishes a voxel oracle from a selector that could motivate
    an adaptive decoder.
\end{itemize}

The study intentionally stops short of claiming an efficient adaptive method.
Measured dynamic compute savings belong to a subsequent method study; here our
goal is to establish whether the empirical premise for such a method is sound.
